\scanpage{162}
where the fiber $E_{\infty}(\unsure{\zeta|\xi})_{\bar s}$ is itself equipped with an abelian extension structure
\[
\tag{342}
\diagram{\begin{tikzcd}[ampersand replacement=\&]
1\arrow[r,dash]\&\underbrace{T_{\infty}(\bar k^{*})}_{T}\arrow[r]\&E_{\infty}(\unsure{\zeta|\xi})_{\bar s}\arrow[r,dash]\&A\arrow[r,dash]\&1.
\end{tikzcd}}
\]

The homomorphism (341) is compatible with this extension structure, when one lets $\Pi_1(S,\bar s)$ act trivially on $A$, and on $T_{\infty}(\bar k^{*})$ \deleted{\unsure{to the} by the}\inserted{(free $\widehat{\mathbb Z}$-module of rank $1$)} by multiplication by the cyclotomic character
\[
\tag{343}
\diagram{\begin{tikzcd}[ampersand replacement=\&]
\Pi_1(S,\bar s)\arrow[r]\&\Pi_1(\mathbb Q,\bar s)\arrow[r,"\chi"]\&\widehat{\mathbb Z}^{*}.
\end{tikzcd}}
\]

Let \deleted{$G$ be the group}
\[
\tag{344}
\begin{gathered}
G_{\unsure{\zeta|\xi}}\subset\operatorname{Aut}(E_{\infty}(\unsure{\zeta|\xi})_{\bar s})\\
\Vert\\
\left\{u\in\operatorname{Aut}\ \middle|
\begin{gathered}
u\text{ respects the extension structure by}\\
\text{inducing }\operatorname{id}_A\ \deleted{\text{of}}\text{ --- and }\inserted{\text{in }T_{\infty}(\bar k^{*})}\text{ a}\\
\text{homothety by an element}\\
\mu\in\widehat{\mathbb Z}^{*}\text{ uniquely determined}
\end{gathered}\right\}
\end{gathered}
\]
\inserted{we thus have a canonical homomorphism}\inserted{$u\mapsto\mu$}
\deleted{a group depends only on $A$}
\[
\diagram{\begin{tikzcd}[ampersand replacement=\&]
G\arrow[r]\&\widehat{\mathbb Z}^{*},
\end{tikzcd}}
\]
fitting into a diagram
\[
\tag{345}
\diagram{\begin{tikzcd}[ampersand replacement=\&]
1\arrow[r]\&\operatorname{Hom}_{\mathbb Z}(A,T_{\infty}(\bar k^{*}))\arrow[r]\&G_{\unsure{\zeta|\xi}}\arrow[r]\&\widehat{\mathbb Z}^{*},
\end{tikzcd}}
\]
\margin{161}

\scanpage{163}
\mspage{729}

where the image \deleted{is} of $G_{\unsure{\zeta|\xi}}\longrightarrow\widehat{\mathbb Z}^{*}$ consists of the $\mu\in\widehat{\mathbb Z}^{*}$ such that \inserted{the homothety,} $\mu_T$ does not change the class of the extension (342) of $A$ by $T$. The homomorphism (341) can thus be interpreted as fitting into a diagram
\[
\tag{346}
\diagram{\begin{tikzcd}[ampersand replacement=\&,column sep=large,row sep=large]
1\arrow[r] \& \operatorname{Hom}_{\mathbb Z}(A,T_\infty(\bar k^{*}))\arrow[r] \& G_{\unsure{\zeta|\xi}}\arrow[r] \& \widehat{\mathbb Z}^{*}\\
\&\& \Pi_1\arrow[u,"\gamma_{\unsure{\zeta|\xi}}"]\arrow[ur,dashed,"{\text{``cyclotomic \unsure{character}''}}"',"{\text{commutative}}" description] \&
\end{tikzcd}}
\]

\margin{N.B.}
If $S$ is \unsure{essentially} of finite type over $\mathbb Q$, then the image of $\Pi_1\longrightarrow\widehat{\mathbb Z}^{*}$ is an open subgroup of $\widehat{\mathbb Z}^{*}$, hence the same is a fortiori true of the image of $G_{\unsure{\zeta|\xi}}\longrightarrow\widehat{\mathbb Z}^{*}$ --- which will even be surjective if the cyclotomic character on $\Pi_1$ is surjective, i.e. if \deleted{\unsure{one}} the only roots of $1$ on $S$ are $\pm1$ (for example if $S=\operatorname{Spec}\mathbb Q\ldots$).

The diagram (346) has \unsure{an} an obvious functoriality property for variable $A$ (say), i.e. for a

\margin{162}

\scanpage{164}

homomorphism of abelian groups
\[
\tag{347}
\diagram{\begin{tikzcd}[ampersand replacement=\&]
A'\arrow[r,"j"] \& A
\end{tikzcd}},
\]
\deleted{the extension}
giving a composite
\[
\tag{348}
\diagram{\begin{tikzcd}[ampersand replacement=\&,column sep=large]
A'\arrow[r,"j"]\arrow[rr,bend right=25,"\unsure{\zeta|\xi}'"'] \& A\arrow[r,"\unsure{\zeta|\xi}"] \& \mathbb G_m(S)
\end{tikzcd}},
\]
the extension $E_\infty(\unsure{\zeta|\xi}')_{\bar s}$ is identified with the inverse image of the extension $E_\infty(\unsure{\zeta|\xi})_{\bar s}$ under (347), whence a homomorphism
\margin{\deleted{\unsure{$G_{\unsure{\zeta|\xi}}$|$\mathbb Q_{\unsure{\zeta|\xi}}$}}}
\[
\diagram{\begin{tikzcd}[ampersand replacement=\&]
G_{\unsure{\zeta|\xi}}\arrow[r] \& G_{\unsure{\zeta|\xi}'}
\end{tikzcd}}
\]
\margin{injective if $j$ is surjective}
fitting into a commutative diagram

\deleted{\unsure{(349)}}
\[
\tag{349}
\diagram{\begin{tikzcd}[ampersand replacement=\&,column sep=large,row sep=large]
1\arrow[r] \& \operatorname{Hom}_{\mathbb Z}(A',T)\arrow[r] \& G_{\unsure{\zeta|\xi}'}\arrow[r] \& \widehat{\mathbb Z}^{*}\\
1\arrow[r] \& \operatorname{Hom}_{\mathbb Z}(A,T)\arrow[u]\arrow[r] \& G_{\unsure{\zeta|\xi}}\arrow[u]\arrow[r] \& \widehat{\mathbb Z}^{*}\arrow[u,"\unsure{1}"]\\
\&\& \Pi_1\arrow[u,"\gamma_{\unsure{\zeta|\xi}}"]\arrow[uu,bend left=35,"\gamma_{\unsure{\zeta|\xi}'}"]\arrow[ur,"{\text{cyclotomic character}}"'] \&
\end{tikzcd}}
\]

where the arrow on the $\operatorname{Hom}_{\mathbb Z}$ is the transpose of (347). This shows that to study \inserted{the structure of} the image of $\gamma_{\unsure{\zeta|\xi}}$ (346), one is reduced, by functoriality, \unsure{to its|to}

\margin{163}

\scanpage{165}
\mspage{730}

study in the case where $\unsure{j|\zeta}$ is \uline{injective}
\[
\tag{350}
\unsure{j|\zeta}:\quad
\diagram{\begin{tikzcd}[ampersand replacement=\&]
A\arrow[r,hook] \& \mathbb G_m(S)
\end{tikzcd}}.
\]

\deleted{I suppose say}
One can moreover, ``divide'' the diagram (346) a little further, by introducing
\[
\tag{351}
\Pi_1^{+}=\Pi_1^{+}(S,\bar s)\overset{\mathrm{def}}=
\operatorname{Ker}\bigl(\Pi_1(S,\bar s)\xrightarrow{\mathrm{cyclot.}}\widehat{\mathbb Z}^{*}\bigr),
\]
and the homomorphism of exact sequences
\[
\tag{352}
\diagram{\begin{tikzcd}[ampersand replacement=\&,column sep=large,row sep=large]
1\arrow[r,dash] \& \Pi_1^{+}\arrow[r,dash]\arrow[d,"\gamma_{\unsure{\zeta|\xi}}^{+}"] \& \Pi_1\arrow[r,"\chi_\Pi"]\arrow[d,"\gamma_{\unsure{\zeta|\xi}}"] \& \widehat{\mathbb Z}^{*}\arrow[d,"\unsure{1}"]\\
1\arrow[r,dash] \& \operatorname{Hom}_{\mathbb Z}(A,T)\arrow[r] \& \deleted{\unsure{H|K}}\inserted{G}_{\unsure{\zeta|\xi}}\arrow[r,dash] \& \widehat{\mathbb Z}^{*}.
\end{tikzcd}}
\]

One can consider that the study of the ``size'' of $\operatorname{Im}(\gamma_{\unsure{\zeta|\xi}}:\Pi_1\longrightarrow G_{\unsure{\zeta|\xi}})$ is \unsure{essentially} reduced to that of the size of $\operatorname{Im}\chi_\Pi$ and of that of $\operatorname{Im}.\gamma_{\unsure{\zeta|\xi}}^{+}$; precisely, if we put
\[
\tag{353}
\operatorname{Im}\gamma_{\unsure{\zeta|\xi}}=\widetilde\Pi_1,\quad
\operatorname{Im}\gamma_{\unsure{\zeta|\xi}}^{+}=\widetilde\Pi_1^{+},\quad
\operatorname{Im}\chi_\Pi=\deleted{\unsure{\widehat{\mathbb Z}^{*}|\widetilde\Gamma}}\inserted{\Gamma}
\]
we will have an exact sequence
\[
\tag{354}
\diagram{\begin{tikzcd}[ampersand replacement=\&]
1\arrow[r] \& \widetilde\Pi_1^{+}\arrow[r] \& \widetilde\Pi_1\arrow[r] \& \Gamma\arrow[r] \& 1
\end{tikzcd}}.
\]
I would like to try to give some

\margin{164}

\scanpage{166}

indications on the size of $\widetilde\Pi_1^{+}\underset{\unsure{\zeta|\gamma}}{\subset}\underbrace{\operatorname{Hom}(A,T)}_{\substack{\simeq\operatorname{Hom}(A,\widehat{\mathbb Z})\\\deleted{\illegible{1}}\inserted{\unsure{\text{noncanonical}^{t}|\text{noncanonical!}}}}}$ in the case where we assume
\[
\tag{355}
\left\{\begin{array}{ll}
1^\circ) & A\quad\mathbb Z\text{-module of finite type}\\
2^\circ) & S\quad\text{essentially}^{\unsure{t|l}}\text{ of finite type over }\mathbb Q\text{ and \uline{normal}}\\
3^\circ) & \unsure{j|\zeta}:A\longrightarrow\mathbb G_m(S)\text{ is injective }\deleted{\unsure{\text{pure}}}.
\end{array}\right.
\]

The hypothesis $2^\circ)$ already implies as we have seen, that $\Gamma\subset\widehat{\mathbb Z}^{*}$ is an open subgroup. The hypothesis that $S$ is normal implies moreover that the groups sought $\widetilde\Pi_1^{+}$ and $\Gamma$ do not change \inserted{\unsure{on$^{h}$|on$^{t}$}} on replacing $S$ by its generic point (we assume for simplicity that $\bar s$ is above the said generic point), so \unsure{suppose} that $S=\operatorname{Spec}K$, where $K$ is a \inserted{field} extension of finite type of $\mathbb Q$. And one can also assume, if one cares to, that $S$ is affine, and smooth over $\unsure{\mathbb Q}\ldots$.

Let us also note, putting
\[
\tag{356}
\check A=\operatorname{Hom}_{\mathbb Z}(A,\mathbb Z),
\]

\margin{165}

\scanpage{167}
\mspage{731}

that $\check A$ is a free $\mathbb Z$-module of finite type (of rank \inserted{$n$} equal to $\operatorname{rg}_{\mathbb Z}A$), and \unsure{that we have} a canonical isomorphism
\[
\tag{357}
\operatorname{Hom}_{\mathbb Z}\left(A,\underbrace{T}_{\substack{\simeq\widehat{\mathbb Z}\\(\unsure{\text{noncanonical}^{t}|\text{noncanonical!}})}}\right)
\simeq
\underbrace{\check A\otimes_{\mathbb Z}T}_{\check A(-1)}
\quad
\left(\begin{gathered}
\simeq\widehat{\mathbb Z}^{n}\\
\unsure{\text{noncanonical}^{t}|\text{noncanonical!}}
\end{gathered}\right)
\]

I would be inclined to think that the image \deleted{\unsure{of}} $\widetilde\Pi_1^{+}$ of $\Pi_1$ in $\operatorname{Hom}_{\mathbb Z}(A,T)\simeq\check A\otimes_{\mathbb Z}T$ is \uline{open}. I even wonder, in the case $S_1$ normal \inserted{$0$-connected} of finite type over $\mathbb Z$, \inserted{dimension $\unsure{2|\geq2}$} and $A=H^0(S_1,\mathbb G_m)$, putting $S=S_{1\mathbb Q}$, whether the homomorphism
\[
\diagram{\begin{tikzcd}[ampersand replacement=\&,column sep=large]
\Pi_1=\Pi_1(S,\bar s)\arrow[r,no head] \& \operatorname{Hom}(A,T)\simeq\check A(-1)
\end{tikzcd}}
\]
is not necessarily surjective??

(More \unsure{generally}, this may be so if $\mathbb G_m(S)/\operatorname{Im}A$ is torsion-free, a case to which one should always be able to reduce, under the hypothesis (355), at the cost of enlarging $A$ to $A'\inserted{\subset H^0(S,\mathbb G_m)}$, \unsure{to|or} an $A'$ such that $A'/A$ is finite\ldots).

\margin{166}

\scanpage{168}

I will return to the case $S=\operatorname{Spec}\mathbb Q$,
then \deleted{\unsure{$\operatorname{Gal}$}}
\[
\mathbb G_m(\unsure{S})=\mathbb G_m(\mathbb Q)=\mathbb Q^*
\simeq\{\pm1\}\times\mathbb Z^{(\mathbb P)}
\]
where $\mathbb P$ is the \unsure{set} of prime \unsure{numbers},
\[
\Pi_1=\operatorname{Gal}(\overline{\mathbb Q}/\mathbb Q)=\Gamma_{\mathbb Q}.
\]

The question is thus to study the homomorphisms
\[
\tag{358}
\unsure{\xi|\zeta}_p:\Gamma_{\mathbb Q}^{+}\longrightarrow\widehat{\mathbb Z}(-1)
\]
associated to the $p\in\mathbb P$, and in particular to answer
the questions

a) are the $\unsure{\xi|\zeta}_p$ surjective, or at least is the image
an open subgroup?

b) are the $\unsure{\xi|\zeta}_p$ linearly independent
over $\widehat{\mathbb Z}$? More precisely, if $p_1\ldots p_n$ are
distinct prime \unsure{numbers}, is the homomorphism
\diagram{\begin{tikzcd}[ampersand replacement=\&,column sep=large]
\Gamma_{\mathbb Q}^{+}\arrow[r,"{(\unsure{\xi|\zeta}_{p_1},\ldots,\unsure{\xi|\zeta}_{p_n})}"']
\&\widehat{\mathbb Z}^{n}(\unsure{-1|1})
\end{tikzcd}}
surjective, or at least is the image of
$\Gamma_{\mathbb Q}^{+}$ an open subgroup?

I will start by looking a little at
question a) (the case at which we
\unsure{were} in the \unsure{context} of Fermat was
the one where $p=2$). For this, we decompose
$\unsure{\xi|\zeta}_p$ into its components
\[
\tag{359}
\unsure{\xi|\zeta}_{p,\deleted{\unsure{\ell}}\ell}:\Gamma_{\mathbb Q}^{+}
\longrightarrow\widehat{\mathbb Z}_{\ell}(\unsure{-1|1})
\qquad(\ell\in\mathbb P)
\]

\margin{167}

\scanpage{169}
\mspage{732}

I would therefore like to prove \inserted{(at least!)} that these are
surjective homomorphisms, or at least \unsure{surjective}
except for a finite \unsure{number} of $\ell$, and with open
image for $\forall\ell$. To study $\unsure{\xi|\zeta}_{p,\ell}$ is
to restrict oneself to the tower of isogenies
\[
\tag{339}
\diagram{\begin{tikzcd}[ampersand replacement=\&,column sep=large]
1\arrow[r,dash]\&\mu_n\arrow[r,dash]\&\mathbb G_m
\arrow[r,"{\lambda\mapsto\lambda^n}"]\&\mathbb G_m\arrow[r,dash]\&1
\end{tikzcd}}
\]
for $n$ a power of $\ell$
\[
\tag{360}
n=\ell^\nu,
\]
and study \deleted{\unsure{the|the}}\inserted{the tower of}
inverse images of $p\in \mathbb G_m(\mathbb Q)$
under its isogenies. It is natural then
to look at ``the'' decomposition group
\margin{\deleted{\unsure{(360)|(359)}}}
\[
\tag{361}
D_\ell\unsure{=|\simeq}\Pi_1(\mathbb Q_\ell,\overline{\mathbb Q}_\ell)
\hookrightarrow\Gamma_{\mathbb Q}
\]
and the cyclotomic \inserted{$\ell$-adic} character induced
\[
\tag{362}
\diagram{\begin{tikzcd}[ampersand replacement=\&,column sep=large]
D_\ell\arrow[r,"\chi_\ell"]\&\mathbb Z_\ell^*
\end{tikzcd}},
\]
or better, the inertia subgroup
\[
\tag{363}
I_\ell\subset D_\ell
\]
and the character induced by $\chi_\ell$ on $I_\ell$.
Indeed we will have
\[
\tag{364}
I_\ell^+\overset{\mathrm{def}}{=}
\operatorname{Ker}(\chi_\ell:I_\ell\longrightarrow\mathbb Z_\ell^*)
\subset\Gamma_{\mathbb Q}^+
\]
for the $\chi_{\ell'}$ $(\ell'\ne\ell)$ are trivial on $I_\ell$
(they are not trivial however on $D_\ell$!).

\margin{168}

\scanpage{170}

We are therefore going to look at the restriction
\[
\tag{365}
\unsure{\xi|\zeta}_{p,\ell}\mid I_\ell^+:
I_\ell^+\longrightarrow\widehat{\mathbb Z}_\ell(-1)
\]
and one can already ask whether this
homomorphism is surjective, or whether
at least its image is \unsure{open|open}?

\margin{$\unsure{=|-}$}

\margin{169}

\scanpage{171}
\mspage{733}

16°) \uline{Reflections on Fermat's theorem.}

``Fermat's theorem'' can be stated by
saying that for $n\geq3$, the affine Fermat
curve over $\mathbb Q$
\[
\tag{366}
\begin{aligned}
U_n\overset{\mathrm{def}}{=}
&\left\{(x_0,x_1,x_\infty)^{\cdot}\in\Gamma\mathbb P^2\ \middle|\
\begin{gathered}
x_0,x_1,x_\infty\ \mathrm{inv.}\\
x_0^n+x_1^n+x_\infty^n=0
\end{gathered}\right\}\\
\simeq{}&\left\{(x_0,\unsure{x_1|a_1},\unsure{x_\infty|a_\infty})^{\cdot}\in
\Gamma \mathbb G_m^{\deleted{\unsure{3}}\inserted{\{0,1,\infty\}}}/\mathbb G_m
\ \middle|\ x_0^n+x_1^n+x_\infty^n=0\right\}
\end{aligned}
\]
has no \deleted{\unsure{solution|solutions}} rational point $/\mathbb Q$.
We had introduced \deleted{\unsure{this}}\inserted{the} ``tower'' of Fermat
curves in §11°) (page 696 and following), and in
14°) there appeared the description of $U_n$ as
the inverse image of
\[
\tag{367}
\begin{array}{l}
U_1\subset\overbrace{\mathbb G_{m_{\mathbb Q}}\otimes_{\mathbb Z}M}^{T}
\simeq \mathbb G_{m_{\mathbb Q}}^{\{0,1,\infty\}}/\mathbb G_{m_{\mathbb Q}}\\[-.25em]
\quad\Vert\\[-.25em]
\left\{(x_0,x_1,x_\infty)^{\cdot}\in\Gamma \mathbb G_m^{\{0,1,\infty\}}/\mathbb G_m
\ \middle|\ x_0+x_1+x_\infty=0\right\}
\end{array}
\]
under the Kummer isogeny of index $n$
\[
\tag{368}
\diagram{\begin{tikzcd}[ampersand replacement=\&,column sep=large]
1\arrow[r]\&\mu_n\otimes_{\mathbb Z}M\arrow[r,dash]\&T
\arrow[r,"{x\mapsto x^n}"]\&T\arrow[r]\&1
\end{tikzcd}}.
\]

Let us calculate the genus \inserted{$g_n$} of $\widehat U_n$, which is a
covering of degree $n^2$ of
$\widehat U_1\underset{\unsure{*|\ne}}{\simeq}
\widehat{\unsure{\mathcal U|\mathcal M}\deleted{\unsure{\mathbb Q|0,3}}}_{0,3}$,
\inserted{with}\deleted{\unsure{totally}}
ramification\deleted{in}\inserted{of degree $n$ \unsure{at}}
the points \inserted{above} $0,1,\infty$, and
étale elsewhere. The Hurwitz formula

\margin{170}